% ============================================================
%  参考答案与提示
% ============================================================

\newpage
\section*{参考答案与提示}

\noindent{\heiti 练习 1.1}\quad C

\smallskip
\noindent{\heiti 练习 1.2}\quad $B=\set{2,3,4}$．

\smallskip
\noindent{\heiti 练习 1.3}\quad 由 $3\in M$，分三种情况讨论：

（1）若 $a=3$，则 $M=\set{1,3,6}$，三个元素互不相同，符合题意；

（2）若 $a^{2}-a=3$，即 $a^{2}-a-3=0$，解得 $a=\dfrac{1\pm\sqrt{13}}{2}$．此时 $M=\set{1,\,a,\,3}$，且 $a\ne 1$、$a\ne 3$，符合题意；

（3）若 $1=3$，不成立．

综上，$a=3$ 或 $a=\dfrac{1\pm\sqrt{13}}{2}$．

\smallskip
\noindent{\heiti 练习 1.4}\quad 分两种情况：

（1）当 $a=0$ 时，方程化为 $2x+1=0$，只有一个解 $x=-\dfrac{1}{2}$，此时 $A=\set{-\dfrac{1}{2}}$，符合题意；

（2）当 $a\ne 0$ 时，方程为一元二次方程，$A$ 中只有一个元素等价于 $\Delta=4-4a=0$，解得 $a=1$，此时 $A=\set{-1}$．

综上，$a=0$ 或 $a=1$．
